% Generated draft: requires mathematical and source review.
\documentclass[11pt]{article}
\usepackage[margin=25mm]{geometry}
\usepackage{fontspec}
\setmainfont{Latin Modern Roman}
\usepackage{amsmath,amssymb,mathtools,unicode-math}
\setmathfont{Latin Modern Math}
\usepackage{enumitem,needspace,xurl,hyperref}
\setlength{\parindent}{0pt}
\setlength{\parskip}{5pt}
\setlength{\emergencystretch}{3em}
\newcommand{\N}{\mathbb{N}}
\newcommand{\Z}{\mathbb{Z}}
\newcommand{\Q}{\mathbb{Q}}
\newcommand{\R}{\mathbb{R}}
\newcommand{\C}{\mathbb{C}}
\newcommand{\F}{\mathbb{F}}
\newcommand{\A}{\mathbb{A}}
\newcommand{\PP}{\mathbb{P}}
\newcommand{\E}{\mathbb{E}}
\newcommand{\eps}{\varepsilon}
\newcommand{\dd}{\,\mathrm{d}}
\begin{document}
\section*{UnsolvedMath ID: 30001091}
\textbf{Source problem code:} catalogue entry
\subsection{Definitions and mathematical statement}
For a transcendental entire map \(f\), its Fatou set is the normality locus of the iterates. A wandering domain is a Fatou component whose successive forward Fatou components are distinct. Limit functions of \(\{f^n|_U\}\) mean subsequential locally uniform limits in the intended, usually spherical, topology. The postsingular set is conventionally
\[
P(f)=\overline{\bigcup_{m\geq0}f^m(S(f))},
\]
where \(S(f)\) is the closed set of finite critical and asymptotic values. Nonrecurrence of a critical point ordinarily means that it does not belong to the omega-limit set of its forward orbit; the source's full nonrecurrence hypotheses are not supplied.

The record bundles these questions:
\begin{enumerate}
\item If \(f\) has a wandering domain \(U\), can its limit functions there be bounded? It is unspecified whether bounded means individually bounded, finite rather than infinite, or uniformly bounded as a collection.
\item If \(P(f)\) is bounded, can \(f\) have wandering domains? Ask also for the stated variant in which all critical points are nonrecurrent; no assumption excluding asymptotic values is added.
\item For \(f_{\lambda,k}(z)=\lambda z e^z+k\), with \(\lambda=e^{2\pi i\theta}\) and the required nonrecurrence assumptions, can \(f_{\lambda,k}\) have a Siegel disk, that is, a maximal domain on which an iterate is conformally an irrational rotation? The assumptions are referred to but not actually stated, and the parameter ranges are incomplete.
\item Is the complement of the hyperbolicity locus precompact in \(\operatorname{Rat}_d\) modulo Möbius conjugacy? Here \(\operatorname{Rat}_d\) denotes degree-\(d\) rational maps; under the usual degree-\(d\geq2\) convention hyperbolicity means that every critical orbit converges to an attracting cycle. Precompactness asks for compact closure in the intended moduli space, not automatically in an added compactification.
\item For \(f_\lambda(z)=\lambda e^z\), characterize parameters for which the nonescaping set is connected, and separately those for which that set together with \(\{\infty\}\) is connected. Under the usual convention the escaping set is \(I(f)=\{z:|f^m(z)|\to\infty\}\), and nonescaping is its complement, not necessarily the bounded-orbit set. The second topology is that of the Riemann sphere.
\item Can Thurston's algorithm be implemented in
\[
O\bigl((\deg(f)+|\Omega_f|)^n\bigr)
\]
iterations, producing either a point sufficiently close to a fixed point of the pullback map or a Thurston obstruction? The input class, \(\Omega_f\), exponent parameter \(n\), closeness tolerance, metric, and meaning of algorithmic cost are unspecified. The printed iteration bound is retained without replacing it by a newly chosen complexity model.
\end{enumerate}
The bundle is not one fully defined proposition; each clause retains its own missing hypotheses.
\subsection{Short English statement}
The source groups questions about wandering domains, bounded postsingular sets, Siegel disks in an explicit entire family, rational-map moduli, exponential connectivity and the complexity of Thurston's algorithm. Every clause is retained with its source-specific gaps.
\subsection{Sources}
\begin{itemize}
\item Local UnsolvedMath registry, permanent ID \(30001091\), source code \texttt{catalogue entry}, ``Wandering Domains, Hyperbolicity, and Thurston Algorithms.'' Registry curation: Ulam AI / UnsolvedMath Contributors, CC BY 4.0; underlying sources retain their supplied provenance.
\item \emph{Trends and Developments in Complex Dynamics} (2009), Oberwolfach Reports, pp.\ 2703--2762, DOI \url{https://doi.org/10.4171/owr/2008/48}.
\item Supplied literature triage, \texttt{literature\_checked\_at}: 2026-08-21: labeled open; says the entry is a bundle of transcendental and rational dynamics questions rather than a single resolved assertion, and the listed parameter-space and algorithmic questions remain an open research bundle. This is supplied dated curation, not independently verified, and gives no precise per-clause status.
\end{itemize}
\end{document}
